% Part: intuitionistic-logic
% Chapter: tableaux
% Section: rules

\documentclass[../../../include/open-logic-section]{subfiles}

\begin{document}

\olfileid{int}{tab}{rul}

\olsection{স্বজ্ঞাবাদী যুক্তিবিদ্যার বিধি}

$\land$ ও $\lor$-এর বিধিগুলি সাধারণ বচনমূলক চিহ্নিত
!!{tableau}s-এর বিধির মতোই, কেবল উপসর্গ যুক্ত হয়। প্রতিটি
ক্ষেত্রে চিহ্নিত !!{formula}~$\sFmla{S}{!A}[\sigma]$-তে
বিধি প্রয়োগ করলে নতুন !!{formula}s পাওয়া যায়, যেগুলির
উপসর্গও~$\sigma$। এটি স্বজ্ঞাগতভাবেই স্পষ্ট: যেমন,
$!A \land !B$ যদি~$\sigma$-নামক জগতে সত্য হয়, তবে
$!A$ ও $!B$-ও~$\sigma$-তে সত্য।
% BN-SRC-410: persistence does not rule out their truth at later worlds.
এই বিধিতে একই জগৎ সম্পর্কে সিদ্ধান্ত টানা হয়; পরবর্তী জগতেও
সত্যতা থাকতে পারে। $\land$ ও $\lor$-এর বিধিগুলি
\olref{tab:prop-rules}-এ দেওয়া হয়েছে।

\begin{table}
  \[\def\arraystretch{3}\begin{array}{|c|c|}
    \hline
    \AxiomC{\sFmla{\True}{!A \land !B}[\sigma]}
    \RightLabel{\TRule{\True}{\land}}
    \UnaryInfC{\sFmla{\True}{!A}[\sigma]}
    \noLine
    \UnaryInfC{\sFmla{\True}{!B}[\sigma]}
    \DisplayProof
    &
    \AxiomC{\sFmla{\False}{!A \land !B}[\sigma]}
    \RightLabel{\TRule{\False}{\land}}
    \UnaryInfC{$\sFmla{\False}{!A}[\sigma] \quad \mid \quad
      \sFmla{\False}{!B}[\sigma]$}
    \DisplayProof
    \\[2ex]
    \hline
    \AxiomC{\sFmla{\True}{!A \lor !B}[\sigma]}
    \RightLabel{\TRule{\True}{\lor}}
    \UnaryInfC{$\sFmla{\True}{!A}[\sigma] \quad \mid \quad
      \sFmla{\True}{!B}[\sigma]$}
    \DisplayProof
    &
    \AxiomC{\sFmla{\False}{!A \lor !B}[\sigma]}
    \RightLabel{\TRule{\False}{\lor}}
    \UnaryInfC{\sFmla{\False}{!A}[\sigma]}
    \noLine
    \UnaryInfC{\sFmla{\False}{!B}[\sigma]}
    \DisplayProof
    \\[2ex]
    \hline
  \end{array}\]
  \caption{উপসর্গযুক্ত !!{tableau}-তে $\land$ ও $\lor$-এর বিধি}
  \ollabel{tab:prop-rules}
\end{table}

বন্ধ হওয়ার শর্তটি সাধারণ !!{tableau}s-এর মতো, তবে এখানে
কেবল !!{formula}s নয়, উপসর্গও মিলতে হবে। আসলে শর্তটি
আরও কিছুটা শিথিল করা যায়। স্বজ্ঞাবাদী মডেলে
!!{formula}s একবার সত্য হলে সত্যই থাকে। তাই~$!A$
$\sigma$-তে সত্য, অথচ প্রবেশযোগ্য উপসর্গ~$\sigma.{*}$-তে
মিথ্যা হওয়া অসম্ভব। সুতরাং কোনো~$\sigma$ ও
!!{formula}~$!A$-র জন্য একটি শাখায় উভয়ই থাকলে শাখাটি বন্ধ:
\[
\sFmla{\True}{!A}[\sigma] \quad\text{and}\quad \sFmla{\False}{!A}[\sigma.{*}]
\]
লক্ষ করুন, চিহ্নগুলি উল্টো হলে—অর্থাৎ শাখায়
\[
\sFmla{\False}{!A}[\sigma] \quad\text{and}\quad \sFmla{\True}{!A}[\sigma.{*}]
\]
থাকলে—শাখাটি কেবল $*$ শূন্য ক্রম হলে বন্ধ।

এ ছাড়া, শাখায়~$\sFmla{\True}{\bot}[\sigma]$ থাকলেও সেটি বন্ধ।

অনুমান স্থাপনের বিধিও সাধারণ !!{tableau}s-এর মতো;
শুধু অনুমানগুলির উপসর্গ সব সময়~$1$। (কোন উপসর্গ নেওয়া
হয়, তা গুরুত্বপূর্ণ নয়—সব অনুমানে একইটি হলেই চলে।)
উদাহরণস্বরূপ,
\[
!B_1, \dots, !B_n \Proves !A
\]
তখনই বলি, যখন নিচের অনুমানগুলি দিয়ে একটি বন্ধ ট্যাবলো হয়:
\[
\sFmla{\True}{!B_1}[1], \dots, \sFmla{\True}{!B_n}[1],
\sFmla{\False}{!A}[1].
\]

শর্তসূচক~$\lif$-এর বিধি ধ্রুপদি ও মোডাল ক্ষেত্র থেকে আলাদা।
% BN-SRC-411: true implication branches on false antecedent or true consequent; sign precedes connective in TRule.
$\TRule{\True}{\lif}$ বিধি, $\sFmla{\True}{!A \lif !B}[\sigma]$
থাকা শাখা থেকে দুটি পৃথক শাখা তৈরি করে: একটিতে
$\sFmla{\False}{!A}[\sigma.{*}]$, অন্যটিতে
$\sFmla{\True}{!B}[\sigma.{*}]$। বিধিটি কেবল সেই
উপসর্গ~$\sigma.{*}$-র জন্য প্রয়োগ করা যায়, যা প্রয়োগের
শাখায় \emph{আগেই} আছে। এমন উপসর্গকে ওই শাখায়
``ব্যবহৃত'' বলব। ($\sigma.{*}$-র মধ্যে $\sigma$ নিজেও
আছে; তাই পৃথক শাখায় উপসর্গযুক্ত চিহ্নিত সূত্র
$\sFmla{\False}{!A}[\sigma]$ ও
$\sFmla{\True}{!B}[\sigma]$ যোগ করে বিধিটি সব সময়
প্রয়োগ করা যায়।)

% BN-SRC-412: TRule takes the sign before the connective here as in the displayed rules.
$\TRule{\False}{\lif}$ বিধি,
$\sFmla{\False}{!A \lif !B}[\sigma]$ থাকা শাখায় একই সঙ্গে
$\sFmla{\True}{!A}[\sigma.n]$ ও
$\sFmla{\False}{!B}[\sigma.n]$ যোগ করে। এখানে
$\sigma.n$ শাখাটিতে নতুন উপসর্গ।

$\lnot$-এর বিধিগুলি একইভাবে সংজ্ঞায়িত হয়
($\lnot !A$-কে $!A \lif \lfalse$ দিয়ে সংজ্ঞায়িত ধরে)।

বিধিগুলি \olref{tab:rules-lif-lnot}-এ দেওয়া হলো।

\begin{table}
  \[\def\arraystretch{3}\begin{array}{|c|c|}
    \hline
    \AxiomC{\sFmla{\True}{\lnot !A}[\sigma]}
    \RightLabel{\TRule{\True}{\lnot}}
    \UnaryInfC{$\sFmla{\False}{!A}[\sigma.{*}]$}
    \DisplayProof
    &
    \AxiomC{\sFmla{\False}{\lnot !A}[\sigma]}
    \RightLabel{\TRule{\False}{\lnot}}
    \UnaryInfC{\sFmla{\True}{!A}[\sigma.n]}
    \DisplayProof
    \\[1ex]
    \text{$\sigma.{*}$ is used} & \text{$\sigma.n$ is new}\\
    \hline
    \AxiomC{\sFmla{\True}{!A \lif !B}[\sigma]}
    \RightLabel{\TRule{\True}{\lif}}
    \UnaryInfC{$\sFmla{\False}{!A}[\sigma.{*}] \quad \mid \quad
      \sFmla{\True}{!B}[\sigma.{*}]$}
    \DisplayProof
    &
    \AxiomC{\sFmla{\False}{!A \lif !B}[\sigma]}
    \RightLabel{\TRule{\False}{\lif}}
    \UnaryInfC{\sFmla{\True}{!A}[\sigma.n]}
    \noLine
    \UnaryInfC{\sFmla{\False}{!B}[\sigma.n]}
    \DisplayProof
    \\[1ex]
    \text{$\sigma.{*}$ is used} & \text{$\sigma.n$ is new}\\
    \hline
  \end{array}\]
  \caption{উপসর্গযুক্ত !!{tableau}-তে $\lnot$ ও $\lif$-এর বিধি}
  \ollabel{tab:rules-lif-lnot}
\end{table}

\end{document}
